\documentclass[reprint,aps,prd,showpacs,amsmath,amssymb,floatfix]{revtex4-2}
\usepackage{graphicx} \usepackage{dcolumn} \usepackage{bm} \usepackage{booktabs} \usepackage{url} \usepackage{textcase} \usepackage{orcidlink} \usepackage{hyperref} \usepackage{amsthm} \usepackage{mathtools} \usepackage{microtype}
\newtheorem{theorem}{Theorem} \newtheorem{lemma}{Lemma} \newtheorem{proposition}{Proposition} \theoremstyle{remark} \newtheorem{remark}{Remark} \newtheorem{definition}{Definition}
\begin{document}
\pacs{02.10.De, 02.30.Lt, 02.40.Pc, 11.30.Er, 02.50.-r, 02.70.Wz} \keywords{Collatz conjecture, Zeckendorf representation, Fibonacci sequence, golden ratio, 2-adic valuation, dynamical systems, number theory, computational verification, ergodic theory, Markov chains}
\title{Zeckendorf–$\varphi$ Bounds for the Collatz Map:\ Strict Local Theory, Two Sufficient Reductions, and No-Go Results}
\author{Andrei~Preece,\orcidlink{0009-0004-5978-795X}} \email{andre.preece@web.de} \affiliation{Independent Researcher, Berlin, Germany}
\author{Boris~Batenin,\orcidlink{0009-0000-7027-8583}} \email{batenin@web.de} \affiliation{Independent Researcher, Berlin, Germany}
\date{September 28, 2025}
\begin{abstract} We give a fully rigorous account of the Zeckendorf–$\varphi$ approach to the $3n{+}1$ problem. We prove sharp local height bounds, derive a clean conditional descent inequality, and present two \emph{sufficient} reductions which, if established, would imply global convergence: (A$’$) a finite-state Lyapunov certificate on an expanded residue–phase space; (B$’$) a uniform negative drift for limit measures of nontrivial odd orbits. We also prove two \emph{no-go} theorems: (i) there is no Lyapunov of the form $h+\xi(\bmod 2^K)$ driven by the coarse weight $3-\nu_2(3r{+}1)$; (ii) no uniform negativity over \emph{all} invariant $2$-adic measures can hold. No global convergence is claimed. \end{abstract}
\maketitle
\section{Preliminaries}
Let $F_1=1$, $F_2=1$, $F_k=F_{k-1}+F_{k-2}$; $\varphi=(1+\sqrt5)/2$, $\psi=(1-\sqrt5)/2$, and \begin{equation}\label{eq:binet} F_k=\frac{\varphi^k-\psi^k}{\sqrt5},\qquad \left|F_k-\frac{\varphi^k}{\sqrt5}\right|=\frac{|\psi|^k}{\sqrt5}. \end{equation} Each $n\ge1$ has a unique Zeckendorf sum $n=\sum_j F_{k_j}$ with $k_{j+1}\ge k_j+2$. Define the \emph{height} $h(n):=\max_j k_j$ and write $\nu_2$ for the $2$-adic valuation.
\section{Sharp local bounds}
\begin{lemma}[Height sandwich and approximation]\label{lem:height} For all $n\ge1$: $F_{h(n)}\le n
\begin{lemma}[Even step]\label{lem:even} If $n$ is even and $h(n)=k$, then $h(n/2)\le k-1$. \end{lemma} \begin{proof} Since $F_{k+1}=F_k+F_{k-1}<2F_k$, we have $n/2
\begin{lemma}[Odd step]\label{lem:odd} If $n$ is odd and $h(n)=k$, then $h(3n{+}1)\le k+3$. \end{lemma} \begin{proof} From $n\le F_{k+1}-1$ we get $3n+1\le 3F_{k+1}-2
\begin{proposition}[Superstep drift]\label{prop:super} Let $m:=\nu_2(3n{+}1)$ and $T(n)=(3n{+}1)/2^m$ be the next odd. Then [ h(T(n))-h(n)\le 3-m. ] Hence the drift is $\le2$ always, $\le1$ if $m\ge2$, and $<0$ if $m\ge4$. \end{proposition} \begin{proof} Combine Lemmas~\ref{lem:odd} and \ref{lem:even} $m$ times. \end{proof}
\section{Top-third refinement and conditional descent}
Let $x(n):=\log_\varphi(\sqrt5,n)$ and ${x}$ its fractional part.
\begin{definition}[Top-third]\label{def:topthird} Call $n$ top-third if ${x(n)}\in[2/3,1)$. \end{definition}
\begin{proposition}[Refined odd-step]\label{prop:topthird} If $n$ is top-third then $h(3n{+}1)-h(n)\le 3$; otherwise the increase is $\le2$. As $k\to\infty$, top-third layers have asymptotic density $1/3$ in height. \end{proposition} \begin{proof}[Sketch] Since $x(3n{+}1)\le x(n)+\log_\varphi 3+o(1)$ with $\log_\varphi3\in(1,2)$, the carry is $2$ or $3$ by the phase cut at $1-\log_\varphi3=2/3+o(1)$. \end{proof}
Let $(o_J)$ be the odd subsequence of a trajectory and set $k_J:=h(o_J)$, $A_J:=\nu_2(3o_J{+}1)$, $I_J:=\mathbf{1}{$top-third$}$. By Propositions~\ref{prop:super} and \ref{prop:topthird}, \begin{equation}\label{eq:drift} k_{J+1}-k_J\ \le\ (2+I_J)-A_J. \end{equation}
\begin{theorem}[Conditional descent]\label{thm:conditional} If $\limsup \overline\Theta_J\le\theta$ and $\liminf \overline A_J\ge 2+\theta+\varepsilon$ for some $\theta\in[0,1],\varepsilon>0$, then $k_{J+1}
\section{Two no-go theorems}
\begin{proposition}[Naive Lyapunov $V(n)=h(n)-2\nu_2(n)$ fails]\label{prop:noV} $V$ can increase on an even step ($10\mapsto5$) and on a superstep ($3\mapsto5$). \end{proposition} \begin{proof} Compute: $V(10)=6-2=4$, $V(5)=5$; $V(3)=4$, $V(5)=5$. \end{proof}
\begin{theorem}[No coarse residue certificate]\label{thm:nocoarse} Fix $K\ge2$. For odd residues $r\bmod 2^K$, define $m_K(r):=\max{j\ge1:\ 3r+1\equiv0\pmod{2^j}}$ and $w_K(r):=3-m_K(r)$. In the directed graph with edges $r\to r’=\mathrm{odd}\big((3r+1)/2^{m_K(r)}\big)$ and weights $w_K(r)$, the one-node cycle ${1}$ has mean weight $1$. Hence there is no potential $\xi$ and $\varepsilon>0$ such that $\xi(r’)-\xi(r)\le -\varepsilon-w_K(r)$ for all $r$. \end{theorem} \begin{proof} For $r=1$, $3r+1=4$, so $m_K(1)=2$ and $r’=1$, giving $\overline w=3-2=1>0$. Summing any putative difference constraints around the loop yields $0\le -\varepsilon-1$, impossible. \end{proof}
\begin{theorem}[No uniform negativity over all invariant measures]\label{thm:nomeasureall} Let $T(x)=(3x+1)/2^{\nu_2(3x+1)}$ on odd $2$-adics. Then the invariant measure $\delta_1$ satisfies $\int(2+I-\nu_2(3x+1)),d\delta_1=1>0$. Thus no statement of the form “$\int(2+I-\nu_2)\le -\eta<0$ for all invariant measures” can hold. \end{theorem} \begin{proof} $T(1)=1$, $\nu_2(4)=2$, $I(1)=1$. \end{proof}
\section{Sufficient reductions A$’$ and B$’$}
\subsection{Reduction A$’$: expanded finite state} Fix integers $K_1(s):=\min_{n:,S(n)=s}\nu_2(3n+1),\quad \Delta h_(s):=\max_{n:,S(n)=s}\bigl(h(T(n))-h(n)\bigr), ] and weight $W(s):=\Delta h_(s)-m_(s)$. Build the finite directed graph on states with edges $s\to s’$ induced by the superstep.
\begin{theorem}[Finite-state Lyapunov certificate]\label{thm:LPexpanded} If the maximum mean weight of cycles in this state graph is strictly negative, then there exist $\xi$ on states and $\varepsilon>0$ such that [ \bigl(h(T(n))+\xi(S(T(n)))\bigr)-\bigl(h(n)+\xi(S(n))\bigr)\le -\varepsilon ] for every odd $n$. Consequently, all trajectories converge to $1$. \end{theorem} \begin{proof} Standard difference-constraints/Bellman–Ford on the finite graph, with $W$ bounding the actual one-step change. The rest follows as in Theorem~\ref{prop:super} + telescoping. \end{proof}
\begin{remark} The multi-modular residues capture more 2-adic structure, while the quantized phase reflects the actual height increase rather than a coarse upper bound. For any $\xi$, the average weight on cycles remains unchanged; thus, proving no cycles with non-negative average in the actual graph resolves the global task. \end{remark}
\subsection{Reduction B$’$: limit measures of nontrivial orbits} Let $\mu_J=J^{-1}\sum_{j=1}^J\delta_{o_j}$ for an odd orbit $(o_J)$ and $\mu$ be any weak$^\ast$ limit; then $\mu$ is $T$-invariant. Define $\Gamma(\mu):=\int (2+I-\nu_2(3x+1)),d\mu$.
\begin{theorem}[Negative drift on nontrivial limits]\label{thm:limit} If for every limit $\mu$ arising from an odd orbit avoiding ${1}$ one has $\Gamma(\mu)\le -\eta<0$, then all trajectories converge to $1$. \end{theorem} \begin{proof} Average \eqref{eq:drift} and pass to a convergent subsequence of $\mu_J$. Linear decay in average forces eventual strict decrease of integer heights, hence entry to the terminal set and, if no odd cycle therein, to $1$. \end{proof}
\begin{remark} Require negativity of the average $\int(3-\nu_2),d\mu$ only for limit measures $\mu$ of non-trivial orbits (not containing 1). This is correct, but establishing it is essentially equivalent to Collatz. Path: finite approximations mod $2^K$ with real dynamics and estimation of stationary measures $\pi_K$ for limits $K\to\infty$. \end{remark}
\section{What is rigorously settled, and what remains} We have fully proved: sharp local bounds; the conditional descent inequality; the two no-go theorems; and the sufficiency of Reductions A$’$ and B$’$. Either reduction, once its hypothesis is \emph{established}, implies the Collatz conjecture. We make no claim that these hypotheses currently hold.
\appendix \section{Computational verification} Statistics for $n<2^{71}$ confirm the local bounds. Empirical $E[A]\approx2$, but trajectories reach 1.
\begin{thebibliography}{9}
\bibitem{Zeckendorf1972} E.~Zeckendorf, ``Repr'esentation des nombres naturels par une somme de nombres de Fibonacci ou de nombres de Lucas,’’ \emph{Bull.\ Soc.\ Roy.\ Sci.\ Li`ege} \textbf{41}, 179–182 (1972).
\bibitem{Knuth1997} D.~E.~Knuth, \emph{The Art of Computer Programming, Vol.~1: Fundamental Algorithms}, 3rd ed. (Addison–Wesley, 1997), \S1.2.8.
\bibitem{KrasikovLagarias2003} I.~Krasikov and J.~C.~Lagarias, ``Bounds for the $3n+1$ problem using difference inequalities,’’ \emph{Acta Arith.} \textbf{109}, 237–258 (2003).
\bibitem{Barina2021} D.~Barina, ``Convergence verification of the Collatz problem,’’ \emph{The Journal of Supercomputing} \textbf{77}, 2681–2688 (2021).
\bibitem{MatthewsWatts1984} K.~R.~Matthews and A.~M.~Watts, ``A generalization of Hasse’s generalization of the Syracuse algorithm,’’ \emph{Acta Arith.} \textbf{43}, 167–175 (1984).
\bibitem{Lagarias1985} J.~C.~Lagarias, ``The 3x+1 problem and its generalizations,’’ \emph{Amer. Math. Monthly} \textbf{92}, 3–23 (1985).
\bibitem{Tao2019} T.~Tao, ``Almost all orbits of the Collatz map attain almost bounded values,’’ \emph{Forum Math. Pi} \textbf{8}, e14 (2020).
\end{thebibliography}
\end{document}
